\chapter{Quality and end-to-end coverage}
\label{ch:qualita-e2e-en}

\lettrine[lines=3]{T}{est coverage} for a project like piTantum
is not a dashboard statistic: it is a parallel manual that
narrates, from the browser's perspective, what the user can
actually do. The frontend Cypress suite, today thirty-four specs
strong, grew by accretion between March and May MMXXVI, until it
covered every list page, every CRUD flow, every dropdown of the
\texttt{/optimize} page and the entire rebirth of
\texttt{/schedule}.

\section{Smoke + workflow}

For each first-class page the suite distinguishes two levels:

\begin{description}
  \item[\textsc{smoke}] (\texttt{*\_smoke.cy.ts}). Loads the
    page, checks that the title appears, that the list is not
    empty (or is a legitimate empty state), that the main
    buttons are clickable. Runs in seconds and acts as a
    regression sentinel.
  \item[\textsc{workflow}] (\texttt{*\_workflow.cy.ts}). Walks
    through a full CRUD path: creates an entity, edits it,
    duplicates it, deletes it, verifies that list and detail
    are coherent at each step. Slower (tens of seconds), more
    brittle, but where serious breakages surface.
\end{description}

\section{Coverage map}

\begin{statsbox}[title={Cypress suite map (May MMXXVI)}]
\begin{tabular}{ll}
  \textsc{master data} &
    \texttt{teachers}, \texttt{subjects}, \texttt{classes},
    \texttt{classrooms}, \texttt{plessi}, \texttt{curricula},
    \texttt{students}, \texttt{groups} \\
  \textsc{assignments} &
    \texttt{assignments\_crud\_workflow},
    \texttt{assignments\_bulk}, \texttt{assignments\_lock} \\
  \textsc{co-teaching} & \texttt{coteaching\_crud\_workflow} \\
  \textsc{constraints} &
    \texttt{constraints\_smoke}, \texttt{constraints\_workflow}
    (DSL editor + 4-step wizard + bulk delete) \\
  \textsc{optimize} &
    \texttt{optimize\_dropdowns},
    \texttt{optimize\_phase\_b\_dropdowns},
    \texttt{optimize\_advanced},
    \texttt{optimize\_launch\_solver} \\
  \textsc{schedule} &
    \texttt{schedule\_calendar},
    \texttt{schedule\_workflow}
    (drag-drop, 4 actions, pool, filter, soft-conflict) \\
  \textsc{logistics} &
    \texttt{logistics\_conflict\_pill} \\
  \textsc{working-hours} &
    internal spec to the working-hours tab, 15/15 green after
    fix \texttt{5c1ba34} \\
  \textsc{monitor} & \texttt{monitor\_smoke},
    \texttt{monitor\_workflow} \\
  \textsc{minor} &
    \texttt{minor\_tabs\_smoke}, \texttt{absences\_smoke},
    \texttt{calendar\_view\_smoke}, \texttt{dashboard\_smoke},
    \texttt{navbar\_completeness} \\
  \textsc{cross-cutting} &
    \texttt{navigation}, \texttt{critical\_workflows},
    \texttt{smoke} \\
\end{tabular}
\end{statsbox}

\section{Writing conventions}

\begin{itemize}
  \item Every interactive element exposes a stable
    \texttt{data-testid}, usually formatted as
    \texttt{<area>-<entity>-<action>}: e.g.
    \texttt{sched-lesson-\{id\}}, \texttt{schedule-pool},
    \texttt{schedule-conflict-modal}. The rule: tests must not
    depend on visible text.
  \item Each spec opens with a \texttt{beforeEach} that visits
    the origin page and waits for a landmark selector (the page
    heading). No arbitrary \texttt{cy.wait(ms)}: wait on
    elements.
  \item Test data uses \texttt{cy.intercept} where it must, but
    the canonical setup is a \texttt{db.sqlite3} pre-seeded via
    \texttt{POST /api/import/profile} before the suite.
\end{itemize}

\section{Lessons from the field}

\begin{erroricomunibox}
\begin{itemize}
  \item Commit \texttt{5c1ba34} reminds us that a
    \emph{namespace import} (\texttt{import * as api}) can be
    \emph{shadowed} by a homonymous local variable:
    \texttt{api.get(\dots)} fails at runtime but TypeScript
    sees no error. Rule of thumb: prefer named imports
    (\texttt{import \{ get, post \} from \dots}) for API
    clients.
  \item Selectors based on localized text are fragile: a
    Italian/English translation breaks twenty specs. Every new
    component added in the April-May MMXXVI cycle exposes
    \texttt{data-testid} from its first commit.
  \item A drag-drop tested in Cypress is not equivalent to a
    real drag-drop: modern browsers emit
    \texttt{dragstart}/\texttt{drop} events with a special
    \texttt{DataTransfer} payload that Cypress does not
    faithfully reproduce. The \texttt{schedule\_workflow} spec
    uses custom helpers (\texttt{trigger('dragstart')},
    \texttt{trigger('drop')} with manual payloads) to simulate
    the interaction: keep this in mind when adding new
    gestures.
\end{itemize}
\end{erroricomunibox}

\section{Running the suite}

\begin{lstlisting}[basicstyle=\ttfamily\footnotesize]
cd webui/frontend
npm run cy:open      # interactive GUI (debug single spec)
npm run cy:run       # full headless (CI / pre-merge)
npm run cy:run -- --spec "cypress/e2e/schedule_*.cy.ts"
\end{lstlisting}

The file \texttt{webui/frontend/E2E\_README.md} describes the
full orchestration pipeline, including the FastAPI backend
restart between suites to guarantee state isolation.
